%! AP Statistics — Unit 1, Lesson 1.7 — Homework ANSWER KEY
\documentclass[10pt]{article}
\usepackage{apstats-article}
\usepackage{apstats-key}

%=============================================================================
\begin{document}

\pageheader{Unit 1, Lesson 1.7}{Homework --- Answer Key}
\begin{tabularx}{\linewidth}{X X X}
Name:\;\ans{Key} & Date:\; & Period:\;
\end{tabularx}

\vspace{0.10in}

% ════════════════════════════════════════════════════════════════════════════
% PART A
% ════════════════════════════════════════════════════════════════════════════
\begin{blurbbox}[Part A \quad Calculate Summary Statistics]{navylight}
\small

Sorted data ($n = 10$): 3, 5, 7, 8, 9, 10, 12, 14, 16, 22

\begin{enumerate}[leftmargin=1.5em, itemsep=5pt]

  \item \ans{$\bar{x} = (3+5+7+8+9+10+12+14+16+22)/10 = 106/10 = \mathbf{10.6}$ hours}

  \item \ans{$n=10$ even; 5th and 6th values are 9 and 10.
        Median $= (9+10)/2 = \mathbf{9.5}$ hours.}

  \item \ans{$Q_1 =$ median of 3,5,7,8,9 $= \mathbf{7}$;\quad
        $Q_3 =$ median of 10,12,14,16,22 $= \mathbf{14}$;\quad
        IQR $= 14 - 7 = \mathbf{7}$ hours}

  \item \ans{Lower fence $= 7 - 1.5(7) = -3.5$;\quad
        Upper fence $= 14 + 1.5(7) = 24.5$.\\
        Since $22 < 24.5$, the value 22 is \textbf{not} an outlier.}

  \item \ans{The \textbf{median (9.5 hours)} is a better description of typical phone
        use because the distribution is slightly skewed right (mean 10.6 $>$ median
        9.5), and the median is more resistant to the higher values in the tail.}

\end{enumerate}

\end{blurbbox}

\vspace{0.10in}

% ════════════════════════════════════════════════════════════════════════════
% PART B
% ════════════════════════════════════════════════════════════════════════════
\begin{blurbbox}[Part B \quad Identify an Outlier \& Choose Statistics]{greenacc}
\small

Sorted data ($n = 12$): 45, 62, 65, 68, 70, 72, 74, 75, 77, 80, 82, 88

\begin{enumerate}[leftmargin=1.5em, itemsep=5pt]

  \item \ans{Median $= (72+74)/2 = \mathbf{73}$\\
        $Q_1 =$ median of lower 6 (45,62,65,68,70,72) $= (65+68)/2 = \mathbf{66.5}$\\
        $Q_3 =$ median of upper 6 (74,75,77,80,82,88) $= (77+80)/2 = \mathbf{78.5}$\\
        IQR $= 78.5 - 66.5 = \mathbf{12}$}

  \item \ans{Lower fence $= 66.5 - 1.5(12) = 66.5 - 18 = \mathbf{48.5}$.\\
        Since $45 < 48.5$, the value 45 \textbf{is an outlier}.}

  \item \ans{Mean $(71.5) <$ Median $(73)$ suggests the distribution is
        \textbf{skewed left} --- the low outlier of 45 pulls the mean below the
        typical center.}

  \item \ans{\textbf{Median + IQR}, because the distribution is skewed left and
        has an outlier (45); both are resistant to the extreme score and give
        a more accurate picture of typical performance on the test.}

\end{enumerate}

\end{blurbbox}

\vspace{0.10in}

% ════════════════════════════════════════════════════════════════════════════
% PART C
% ════════════════════════════════════════════════════════════════════════════
\begin{blurbbox}[Part C \quad AP-Style Free Response]{redacc}
\small

\begin{enumerate}[leftmargin=1.5em, itemsep=5pt]

  \item \ans{Because the mean (82) is less than the median (91), the distribution
        is likely \textbf{skewed left} --- there are a few unusually low values
        that pull the mean below the center, while the median resists this effect.}

  \item \ans{The \textbf{median (91)} is more appropriate because the distribution
        is skewed left; the mean is pulled downward by extreme low values and
        does not represent a typical value as well as the resistant median.}

\end{enumerate}

\end{blurbbox}

\end{document}
